\section{风险、收益与投资组合}
\label{sec:09-Convex-Optimization/07-Applications/04}

接手一个跑了很久的配置脚本：输入是六十只标的过去五年的日收益，输出是一组权重。这一版跑出来的结果是第一只 \(+340\%\)，第三只 \(-280\%\)，另有七八只在 \(\pm100\%\) 上下，剩下五十来只几乎为零。求解器的日志很体面，KKT 残差 \(3\times10^{-13}\)，对偶间隙 \(10^{-11}\)。

把估计窗口往后挪一个月重跑，权重换了一批标的，符号也变了，前一版重仓的那只这次被做空。两次的输入差别只是多了二十个交易日、少了二十个交易日。

求解器没有出错，它每次都精确地解出了被交给它的那个问题。出问题的是被交给它的那个问题——更准确地说，是那个问题里被当成已知量的输入。这一篇把这条链条走一遍：先看凸性在这里给了什么，再看它没有给什么。

\subsection{权衡本身可以写成一个二次规划}

设有 \(d\) 项资产，权重向量 \(w\in\R^d\) 表示资金的分配比例。若各资产收益率构成随机向量 \(R\)，期望为 \(\mu=\E[R]\)，协方差为 \(\Sigma=\Cov(R)\)，那么组合收益 \(R^{\trans}w\) 的期望与方差是

\[
\E[R^{\trans}w]=\mu^{\trans}w,\qquad \Var(R^{\trans}w)=w^{\trans}\Sigma w .
\]

第二个式子说明了为什么协方差矩阵必须半正定：方差不能为负，而 \(w\) 可以取任何方向，所以 \(w^{\trans}\Sigma w\ge0\) 对一切 \(w\) 成立，这正是 \(\Sigma\succeq0\) 的定义。这一条不是形式上的约定，它同时是目标函数凸性的全部来源——二次型 \(w^{\trans}\Sigma w\) 凸当且仅当 \(\Sigma\succeq0\)。半正定矩阵全体构成一个凸锥，它的几何见\hyperref[sec:09-Convex-Optimization/02-Convex-Sets/03]{``从方向的集合到半正定锥''}。

要留意的是，半正定是矩阵级别的性质，逐个元素检查查不出来。每只标的的方差为正、每对相关系数落在 \([-1,1]\) 内，这些条件全部满足，矩阵仍可能有负特征值——三个两两相关系数各自合法，凑在一起未必对应任何真实的联合分布。手工调整过、按不同时间窗分块拼接过、或者用不同样本量逐对估计过的相关矩阵，都要重新验一遍最小特征值。

Markowitz 的均值---方差模型有两种等价写法。给定风险厌恶系数 \(\gamma>0\)：

\[
\begin{aligned}
\text{minimize}\quad & \tfrac{\gamma}{2}\,w^{\trans}\Sigma w-\mu^{\trans}w\\
\text{subject to}\quad & \ind^{\trans}w=1,\qquad w\in C .
\end{aligned}
\]

或者给定收益下限 \(r\)，压低风险：

\[
\begin{aligned}
\text{minimize}\quad & w^{\trans}\Sigma w\\
\text{subject to}\quad & \mu^{\trans}w\ge r,\quad \ind^{\trans}w=1,\quad w\in C .
\end{aligned}
\]

这里 \(\ind\) 是全一向量，预算约束 \(\ind^{\trans}w=1\) 说资金全部配出去，\(C\) 收纳其余持仓限制，禁止做空就是 \(w\ge0\)。两种写法都是凸二次规划：目标二次且凸，约束仿射，见\hyperref[sec:09-Convex-Optimization/04-Convex-Problems/02]{``最小二乘、范数与二次规划''}。它们通过 Lagrange 乘子互相对应，但 \(\gamma\) 与 \(r\) 之间没有脱离数据的换算公式，同一个 \(\gamma\) 在不同的 \(\mu,\Sigma\) 下对应不同的 \(r\)。

\subsection{有效前沿是一条曲线，不是一个点}

把 \(r\) 从低到高扫一遍，每个 \(r\) 得到一个最优风险，把 \((\text{风险},\ r)\) 画在平面上，得到的就是有效前沿。

\begin{center}
\begin{tikzpicture}[font=\small,>=stealth,scale=1.25]
  \fill[black!5] plot[domain=0:2.6,samples=70]
    ({0.6+1.5*sqrt(1+((\x-1.1)/0.9)*((\x-1.1)/0.9))},{\x})
    -- (4.8,2.6) -- (4.8,0) -- cycle;
  \draw[black!25,->] (0.9,0) -- (4.85,0);
  \draw[black!25,->] (0.9,-0.1) -- (0.9,3.0);
  \node[font=\footnotesize,black!60] at (3.1,-0.42) {风险 \(\sqrt{w^{\trans}\Sigma w}\)};
  \node[font=\footnotesize,black!60,anchor=west] at (0.85,3.15) {收益 \(\mu^{\trans}w\)};
  \draw[black!35,dashed] plot[domain=0:1.1,samples=40]
    ({0.6+1.5*sqrt(1+((\x-1.1)/0.9)*((\x-1.1)/0.9))},{\x});
  \draw[black!80,thick] plot[domain=1.1:2.6,samples=50]
    ({0.6+1.5*sqrt(1+((\x-1.1)/0.9)*((\x-1.1)/0.9))},{\x});
  \fill[black!80] (2.1,1.1) circle (1.8pt);
  \draw[black!45] (1.78,0.76) -- (2.02,1.03);
  \node[font=\footnotesize,black!70] at (1.5,0.62) {最小方差};
  \fill[black!80] (2.18,1.4) circle (1.8pt);
  \node[font=\footnotesize,black!70] at (1.66,1.45) {\(\gamma\) 大};
  \fill[black!80] (2.97,2.2) circle (1.8pt);
  \node[font=\footnotesize,black!70] at (2.45,2.28) {\(\gamma\) 小};
  \draw[black!60] (2.90,0.60) circle (1.6pt);
  \draw[black!60] (3.60,1.50) circle (1.6pt);
  \draw[black!60] (3.50,2.15) circle (1.6pt);
  \draw[black!60] (4.15,2.45) circle (1.6pt);
  \draw[black!60] (2.75,1.00) circle (1.6pt);
  \draw[black!60] (3.35,0.85) circle (1.6pt);
  \node[font=\footnotesize,black!60] at (3.95,0.42) {单只标的};
\end{tikzpicture}

\emph{阴影是所有可行组合能达到的风险---收益区域，空心点是单只标的；左边界的实线部分是有效前沿，虚线部分被最小方差组合占优。}
\end{center}

图上没有``最优组合''这个东西。实线上的每一点在自己那个风险水平上都无可挑剔：同样风险下没有更高的收益，同样收益下没有更低的风险。它们之间的取舍不在数学里，在委托人的风险偏好里，而 \(\gamma\) 就是这个偏好的数值化身。虚线那一段则确实是坏的：它下方存在收益更高而风险相同的点，任何理由都不足以选它。空心点全都落在前沿右侧，这是分散化的直接图示——把它们混起来，能到达任何单只标的都到不了的左上方。

前沿的形状还解释了一件事：曲线在最小方差组合附近很平。这意味着在那一带，收益目标提高一点，需要付出的风险代价很小；但同样也意味着，最优权重在那一带对输入的变化非常敏感——目标函数在一大片权重上几乎一样好，谁被选中就取决于第几位小数。这正是开头那个脚本的处境。

\subsection{分散靠的是不齐步走，不是数个数}

把二次型展开：

\[
w^{\trans}\Sigma w=\sum_i w_i^2\sigma_i^2+2\sum_{i<j}w_iw_j\Cov(R_i,R_j).
\]

第一项是各自的方差按平方权重加总，第二项是交叉协方差。持有六十只齐步涨跌的标的，交叉项全部为正且接近各自方差的量级，组合方差几乎不比单只低；持有两只弱相关或负相关的标的，交叉项为负，组合方差可以低于两者中的任何一个。分散化是靠协方差的符号做出来的，不是靠标的的数量。

同一个交叉项也是麻烦的来源。相关性越高，\(\Sigma\) 越接近奇异，最优化问题在``此消彼长''的那些方向上就越平。这与前两篇里的现象是同一个：那里是两列共线的特征让 \(X^{\trans}X\) 接近奇异，这里是两只齐步走的标的让 \(\Sigma\) 接近奇异。平谷所在之处，答案就不稳。

\subsection{优化会精确地钻输入误差的空子}

拿最小方差组合算笔账，两只标的就够。约束是 \(w_1+w_2=1\)，最小化 \(w^{\trans}\Sigma w\)，闭式解是

\[
w_1^\star=\frac{\sigma_2^2-\rho\sigma_1\sigma_2}{\sigma_1^2+\sigma_2^2-2\rho\sigma_1\sigma_2},\qquad w_2^\star=1-w_1^\star .
\]

分母正是两只收益之差的方差 \(\Var(R_1-R_2)\)。\(\rho\) 越接近 \(1\)，两只走得越像，这个差的方差越小，分母越小。

代进数字。取 \(\sigma_1=0.20\)、\(\sigma_2=0.22\)、\(\rho=0.99\)：分子是 \(0.0484-0.99\times0.044=0.00484\)，分母是 \(0.0884-0.08712=0.00128\)，于是 \(w^\star=(3.78,\,-2.78)\)——为了压掉一点点方差，模型要求融券卖出近三倍本金的第二只，再把所得全部押到第一只上。现在把 \(\sigma_2\) 的估计从 \(0.22\) 改成 \(0.21\)，一个百分点的波动率估计误差，重算得到 \(w^\star=(2.68,\,-1.68)\)。头寸变了一整倍本金。

机制在分母上：\(\sigma_2\) 的误差进入分子时是一阶的，进入分母时被 \(2\rho\sigma_1\sigma_2\) 这一项几乎抵消，剩下的是两个大数相减的小差。误差没有被平均掉，它被除以了一个接近零的数。这就是\hyperref[sec:08-Numerical-Analysis/01-Floating-Point-and-Error/03]{``条件数衡量问题本身的敏感性''}描述的放大，只不过放大的对象不是舍入误差而是统计估计误差，而放大倍数由 \(\Sigma\) 的条件数给出。

维数升高只会更糟。\(d\) 只标的的协方差矩阵有 \(d(d+1)/2\) 个待估参数，六十只就是一千八百三十个，而五年日收益只有约一千两百五十个观测。样本协方差矩阵此时必定奇异，它最小的那些特征值对应的方向纯粹是采样噪声，见\hyperref[sec:11-Mathematical-Statistics/09-Random-Matrix-Theory/01]{``样本协方差在高维为何会失真''}。而最小方差优化恰恰专门去找方差最小的方向——它会径直走向那些被低估得最厉害的方向，因为在它眼里那里风险最低。预期收益 \(\mu\) 的估计更不可靠，收益率的信噪比低到需要几十年数据才能把均值估准，而优化器会把 \(0.1\%\) 的收益差异当成确凿的事实，据此建立极端头寸。

这个毛病有个准确的名字叫误差最大化：均值---方差优化不是对估计误差不敏感，而是系统性地偏爱那些误差方向。它挑出的组合，恰好是那些``看上去''收益最高、风险最低的组合，而看上去最好通常意味着估计得最离谱。这是\hyperref[sec:09-Convex-Optimization/01-Optimization-Foundations/01]{``把现实选择翻译成优化问题''}里那句话的又一个实例：目标写歪了，求解器会精确地找到那个空子。凸性保证了你找到的是全局最优，它保证不了这个全局最优对应的现实决策是好的。

\subsection{能做的事是把可行域收紧}

对付这件事有两条路，方向相反但常常一起用。

一条是修输入。对 \(\Sigma\) 做收缩估计，把样本协方差与一个结构简单的目标（比如对角阵或单因子模型的隐含协方差）按权重混合，等价于给谱做了一次抬升——这与第一篇里 \(\lambda I\) 给 Hessian 补曲率是同一个动作，只不过发生在输入端而不是目标端。因子模型把 \(d(d+1)/2\) 个参数压到少数几个因子上，也是同类思路。对 \(\mu\) 则做更强的收缩，把点估计往截面均值拉，或者干脆放弃对 \(\mu\) 的估计，只做最小方差组合。

另一条是收紧可行域，让优化器没有空子可钻。单只权重上下界 \(l\le w\le u\)、总杠杆 \(\norm{w}_1\le L\)、相对现有持仓的换手限制 \(\norm{w-w^{\text{old}}}_1\le T\)、行业暴露上限，这些全是仿射或范数球约束，加进去问题仍然是凸的。它们的作用不是表达偏好，而是承认输入不可靠：既然模型算出的 \(340\%\) 头寸建立在一个方差估计的第二位小数上，那就干脆不让它到达那里。禁止做空这条老规矩之所以在实证中表现良好，原因也在这里——\(w\ge0\) 把可行域砍掉一大半，砍掉的恰好是杠杆最容易失控的部分。这些约束在数学上是可行域的边，在统计上起的却是正则化的作用，和 \(\ell_1\) 球限制系数是同一类操作。

再进一步是鲁棒优化：不再假定 \(\mu\) 是已知常量，而是设它落在某个不确定集里，对集合内的最坏情形做优化。若不确定集取成椭球，最坏情形的目标含有一项 \(\norm{\Sigma_\mu^{1/2}w}_2\)，问题变成二阶锥规划，仍然凸，仍然可解。这条路把``我不知道 \(\mu\) 到底是多少''直接写进了模型，而不是让它以极端头寸的形式在输出端爆发。

也有一些现实约束会把凸性彻底破坏。基数约束``最多持有 \(k\) 只''、最小交易单位、固定手续费（交易与否是零一决策），它们引入的是离散结构，问题变成混合整数二次规划，可行域不再是凸集，全局最优不再由局部条件认证。这时的选择是要么上整数规划求解器，要么做凸松弛，而后者必须在报告里写清楚：解出来的是替代问题的全局最优，见\hyperref[sec:09-Convex-Optimization/09-Complexity-and-Approximation/02]{``凸松弛把 NP-hard 问题变成可解问题''}。

\subsection{对偶变量告诉你瓶颈在哪}

凸性在这里还留下一样有用的东西。收益下限约束 \(\mu^{\trans}w\ge r\) 的乘子，等于最优风险对 \(r\) 的变化率：乘子大，说明收益要求再提一点，风险代价陡增，这条约束正卡着脖子；某个行业上限的乘子为零，说明它当前根本没有起作用，放宽它不会改善任何东西。这就是\hyperref[sec:09-Convex-Optimization/05-Duality/04]{``对偶变量是约束的影子价格''}在配置问题上的读法，它把``该往哪里争取放宽''变成一张可以排序的表。

前提照旧要说明白：乘子是当前输入、当前活跃约束集下的局部灵敏度。\(\mu\) 更新一次，活跃集可能整体换掉，乘子跟着变。把某一天算出的影子价格当成长期规律写进制度，是把局部导数当成了全局结论。

\subsection{凡是优化一个估计出来的目标，都有这一课}

这一单元的四篇讲的是同一件事的四个侧面。经验风险最小化把学习写成优化，\(\ell_2\) 正则同时买下唯一性、收敛速率与稳定性；\(\ell_1\) 的尖角把稀疏做成优化结构的产物；最大间隔让少数样本撑起边界；均值---方差把风险偏好写成一个参数。四次都是凸问题，四次都能拿到全局最优与可核验的证书。

四次也都撞上同一堵墙。凸性负责的是``给定这个目标和这个可行域，我找到了全局最优，并且能证明''。它不负责这个目标是不是你真正关心的东西，更不负责目标里那些被当作已知量的数字是从哪里估出来的、带多大误差。开头那个 \(3\times10^{-13}\) 的 KKT 残差是真的，它说明求解器把一个建立在噪声上的问题解得极其精确。

这一课在数据工作里到处都是。用交叉验证挑出的超参数、用历史日志估出的点击率、用少量标注估出的类别先验、用一个模型的输出当作另一个模型的输入——只要下游是一个优化步骤，上游的估计误差就会被有选择地放大，因为优化的职责就是去找极端。防御的手段也总是那几样：收缩输入、收紧可行域、对输入扰动做灵敏度分析、把不确定性写进模型而不是等它从输出端冒出来。

至此，凸优化能给的保证已经交代完整。接下来要面对的是保证本身失效的地方：目标不再凸，梯度只能算出一个带噪的估计，而这恰恰是训练神经网络的日常。\hyperref[ch:078]{``非凸与随机优化：能下降不等于找到全局解''}接着往下走。
